\documentclass[a4paper]{article}
% Español (México) con XeLaTeX: fontspec para fuentes Unicode, polyglossia
% para la división silábica y los nombres del documento en la variante mexicana.
\usepackage{fontspec}
\usepackage{polyglossia}
\setdefaultlanguage[variant=mexican]{spanish}
% Los nombres chinos van como «pinyin (original)»: xeCJK da a los caracteres
% Han su propia fuente; sin ella XeLaTeX los omite con un aviso.
\usepackage{xeCJK}
\setCJKmainfont{FandolSong-Regular.otf}
% polyglossia no traduce los nombres de listings.
\AtBeginDocument{\providecommand{\lstlistingname}{}\renewcommand{\lstlistingname}{Listado}%
  \providecommand{\lstlistlistingname}{}\renewcommand{\lstlistlistingname}{Índice de listados}}
\usepackage{amsmath, amssymb}
\usepackage{graphicx}
\usepackage[margin=2.5cm]{geometry}
\usepackage[most]{tcolorbox}
\usepackage{etoolbox}
\usepackage{listings}
\usepackage{hyperref}
\usepackage{booktabs}
\usepackage{subcaption}
\usepackage{float}

\newtcolorbox{knowledgebox}[1]{enhanced, colback=blue!5!white, colframe=blue!75!black, colbacktitle=blue!75!black, coltitle=white, fonttitle=\bfseries, title={#1}, attach boxed title to top left={yshift=-2mm, xshift=2mm}, boxrule=1pt, sharp corners}
\newtcolorbox{importantbox}[1]{enhanced, colback=yellow!10!white, colframe=yellow!80!black, colbacktitle=yellow!80!black, coltitle=black, fonttitle=\bfseries, title={#1}, sharp corners}
\newtcolorbox{warningbox}[1]{enhanced, colback=red!5!white, colframe=red!75!black, colbacktitle=red!75!black, coltitle=white, fonttitle=\bfseries, title={#1}, sharp corners}

\lstset{language=Python, basicstyle=\ttfamily\small, keywordstyle=\color{blue}, stringstyle=\color{red!60!black}, commentstyle=\color{green!60!black}, breaklines=true, frame=single, numbers=left, numberstyle=\tiny\color{gray}, captionpos=b, extendedchars=true}

\newcommand{\notetitle}{Recuperación de doble capa en el grafo de conocimiento de LightRAG y LangFuse}
\newcommand{\noteauthors}{Elaboradas a partir de materiales públicos del curso}
\newcommand{\notedate}{2025}
\newcommand{\videochannel}{Wudaokou Nash (五道口纳什)}
\newcommand{\videopublishdate}{2025}
\newcommand{\videoduration}{aproximadamente 19 minutos}
\newcommand{\videourl}{https://www.bilibili.com/video/BV1TkmqBKEid}
\newcommand{\videocoverpath}{cover.jpg}
\newcommand{\repourl}{https://github.com/hqhq1025/ai-course-notes}


% Footer with repo info
\usepackage{fancyhdr}
\pagestyle{fancy}
\fancyhf{}
\fancyfoot[C]{\thepage}
\fancyfoot[R]{\footnotesize\color{gray}\href{\repourl}{AI Course Notes}}
\renewcommand{\headrulewidth}{0pt}
\renewcommand{\footrulewidth}{0pt}

\begin{document}
\begin{titlepage}
\centering
{\Large Notas del curso\par}\vspace{1.2cm}
{\huge\bfseries \notetitle\par}\vspace{0.8cm}
{\large \noteauthors\par}\vspace{0.3cm}
{\large \notedate\par}\vspace{1.2cm}
\ifdefempty{\videocoverpath}{}{\includegraphics[width=0.82\textwidth,height=0.45\textheight,keepaspectratio]{\videocoverpath}\par}
\vfill
\begin{tcolorbox}[width=0.9\textwidth, colback=black!2!white, colframe=black!60, sharp corners]
\textbf{Autor o canal del video}: \videochannel\par\textbf{Fecha de publicación}: \videopublishdate\par\textbf{Duración del video}: \videoduration\par
\ifdefempty{\videourl}{}{\textbf{Enlace del video}: \href{\videourl}{\nolinkurl{\videourl}}\par}\par
\textbf{Página del proyecto}: \href{\repourl}{\nolinkurl{\repourl}}\par
\end{tcolorbox}
\end{titlepage}
\tableofcontents\newpage

\section{Cómo familiarizarse rápido con un proyecto de Agent}

\subsection{Síntesis de la metodología}

El autor sintetizó una ruta de tres pasos para comprender y familiarizarse rápido con un proyecto de Agent:

\begin{enumerate}
    \item \textbf{Comprender los conceptos básicos}: leer el artículo o la introducción al framework del proyecto, comprender el diseño del Workflow / Pipeline
    \item \textbf{Comprender el Prompt}: en un proyecto de Agent, gran cantidad de subprocesos son representados por el LLM; comprender el Prompt equivale a comprender la entrada, la salida y la lógica de procesamiento de cada subproceso
    \item \textbf{Usar herramientas de monitoreo}: usar herramientas como LangFuse para hacer trace de todas las llamadas a la API y comprender de forma visual el proceso completo
\end{enumerate}

\begin{knowledgebox}{El Prompt es la clave para comprender un proyecto de Agent}
En el Workflow de un Agent, muchos subprocesos son representados por el LLM. Si se entiende cada subproceso como $F(X; \text{Prompt})$ —donde $F$ es el LLM, $X$ es la Query y Prompt es el System/User Prompt—, comprender el Prompt equivale a comprender la lógica de procesamiento de cada subproceso.
\end{knowledgebox}

\subsection{Resumen de la sección}

Concepto $\rightarrow$ Prompt $\rightarrow$ herramienta de monitoreo: esta es una ruta eficiente para comprender un proyecto de Agent.

\section{Principios de LightRAG}

\subsection{Panorama general}

LightRAG es un artículo académico de EMNLP 2025, además de otro proyecto del autor de Paper2Slides. Es un sistema RAG basado en un \textbf{grafo de conocimiento}, con una capacidad de comprensión semántica más sólida que el RAG vectorial tradicional.

\begin{importantbox}{Características centrales de LightRAG}
\begin{itemize}
    \item Indexación basada en \textbf{Knowledge Graph} (en lugar de un índice puramente vectorial)
    \item \textbf{Recuperación en dos niveles}: Low-Level (local, entidades específicas) + High-Level (global, temas y conceptos)
    \item \textbf{Actualización incremental}: los documentos nuevos no requieren reconstruir todo el Graph; basta con ejecutar el paso de extracción y fusionar
\end{itemize}
\end{importantbox}

\subsection{Proceso de Indexing: construcción del grafo de conocimiento}

Para cada documento se ejecutan los siguientes pasos:

\begin{enumerate}
    \item \textbf{Chunking}: dividir el documento en fragmentos
    \item \textbf{Entity \& Relationship Extraction}: extraer entidades y relaciones de cada Chunk mediante un LLM
    \item \textbf{Profiling}: generar descripciones textuales de las entidades y relaciones mediante un LLM
    \item \textbf{Deduplication}: deduplicar y fusionar en todo el grafo de conocimiento
    \item \textbf{Embedding}: construir un Embedding Vector para la Key de todas las entidades y relaciones
\end{enumerate}

Estructura de datos del grafo de conocimiento:
\begin{itemize}
    \item \textbf{Entity}: nombre, tipo, descripción (generada por Profiling), ID del Chunk original
    \item \textbf{Edge}: Source, Target, palabras clave, descripción, ID del Chunk original
\end{itemize}

\begin{knowledgebox}{LightRAG frente a GraphRAG}
Una ventaja importante de LightRAG frente a GraphRAG es la \textbf{actualización incremental}: para un documento nuevo basta con ejecutar los mismos pasos de extracción y fusionar las entidades y los bordes. En cambio, GraphRAG necesita rehacer la detección de comunidades, lo cual resulta costoso.
\end{knowledgebox}

\subsection{Proceso de Retrieval: recuperación en dos niveles}

\begin{enumerate}
    \item A partir de la Query, el LLM genera \textbf{Low-Level Keywords} (entidades específicas) y \textbf{High-Level Keywords} (temas o conceptos)
    \item Se emparejan los nodos relevantes del grafo de conocimiento según la similitud de Embedding
    \item Se amplía el Subgraph mediante \textbf{vecinos a un salto}
    \item Se obtienen tres tipos de Context: entidades, relaciones, Chunk original
    \item Context + Query $\rightarrow$ el LLM genera la respuesta final
\end{enumerate}

\subsection{Resumen de la sección}

LightRAG logra, mediante la indexación por grafo de conocimiento y la recuperación en dos niveles, una capacidad de comprensión semántica más sólida que el RAG puramente vectorial, además de admitir actualización incremental.

\section{LangFuse: herramienta de monitoreo de API con despliegue local}

\subsection{Introducción a LangFuse}

LangFuse es una herramienta de monitoreo de API de LLM que se puede desplegar localmente, y se usa para hacer trace de todo el proceso de llamadas al modelo.

\subsection{Despliegue y configuración}

\begin{enumerate}
    \item Descargar el archivo de Docker Compose e iniciar el servicio
    \item Acceder localmente a \texttt{localhost:3000}
    \item Crear la organización y el proyecto
    \item Obtener el Public Key, el Secret Key y el Base URL
    \item Configurar la información anterior en \texttt{.env}
\end{enumerate}

LightRAG integra \textbf{de manera nativa} LangFuse: en cuanto detecta las variables de entorno, reemplaza automáticamente el OpenAI Client por el LangFuse Client, y registra en segundo plano los traces completos.

\subsection{Entender el Workflow a través de los Traces}

En el panel de Traces de LangFuse se pueden ver todas las llamadas a la API de LightRAG:

\textbf{Etapa de Insert (Indexing)}:
\begin{enumerate}
    \item Crear el Embedding del documento
    \item Primera Generation: extraer entidades y relaciones (Entity Extraction System/User Prompt)
    \item Segunda Generation: Self Correction (seguir extrayendo el contenido omitido o con formato incorrecto)
    \item Gran cantidad de llamadas de Embedding: crear el Embedding de los nodos y las aristas
\end{enumerate}

\textbf{Etapa de Query (Retrieval)}:
\begin{enumerate}
    \item Generation: extraer palabras clave High-Level y Low-Level (Keywords Extraction Prompt)
    \item Embedding Search: hacer coincidir entidades y relaciones a partir de las palabras clave
    \item Recorrido del grafo: encontrar los nodos vecinos de las entidades relevantes
    \item Generation: generar la respuesta final a partir de Context + Query (RAG Response Prompt)
\end{enumerate}

\begin{warningbox}{Response Cache}
LightRAG tiene habilitado por defecto el Response Cache: si se ejecuta de nuevo la misma Query, no se vuelve a llamar a la API. Hay que tener esto en cuenta al depurar durante el desarrollo.
\end{warningbox}

\subsection{Resumen de la sección}

LangFuse ofrece una capacidad de trace visual poderosa, y es una herramienta útil para entender el mecanismo interno de funcionamiento de un proyecto de Agent.

\section{Análisis de los Prompts centrales de LightRAG}

\subsection{Entity Extraction System Prompt}

Se encarga de extraer entidades y relaciones a partir del Chunk, e incluye ejemplos de Few-Shot.

\subsection{Continue Extraction User Prompt}

A partir del resultado de la primera extracción, realiza Self Correction, es decir, identifica y extrae el contenido omitido o con formato incorrecto.

\subsection{Keywords Extraction Prompt}

Se encarga de extraer palabras clave de High-Level y Low-Level a partir del Query, para usarlas en la búsqueda vectorial posterior.

\subsection{RAG Response Prompt}

A partir del Context recuperado (entidades, relaciones, Chunk) y del Query del usuario, genera la respuesta final.

\subsection{Resumen de la sección}

Los Prompts centrales de LightRAG incluyen: extracción de entidades y relaciones, Self Correction, extracción de palabras clave y generación de la respuesta de RAG; cada uno corresponde a una etapa clave del Workflow.

\section{Estructura de almacenamiento local de LightRAG}

Al ejecutar LightRAG, en el directorio de trabajo se generan los siguientes archivos:

\begin{itemize}
    \item \texttt{graph.graphml}: grafo de conocimiento en formato XML (entidades + relaciones)
    \item \texttt{kv\_store\_*}: almacén clave-valor (pares K-V de Chunk, Entity y Relation)
    \item \texttt{vdb\_*}: base de datos vectorial (usa Nano Vector DB de manera predeterminada)
    \begin{itemize}
        \item \texttt{vdb\_chunks}: Embedding de los Chunk
        \item \texttt{vdb\_entities}: Embedding de las Entity
        \item \texttt{vdb\_relationships}: Embedding de las Relation
    \end{itemize}
\end{itemize}

Flujo de datos durante la recuperación: palabras clave $\rightarrow$ Embedding Search (VDB) $\rightarrow$ recorrido del grafo (GraphML) $\rightarrow$ complemento con el documento original (KV Store) $\rightarrow$ ensamblado del Context $\rightarrow$ generación de la respuesta por el LLM.

\section{Síntesis y ampliación}

\begin{enumerate}
    \item \textbf{Ruta rápida para iniciarse en un proyecto de Agent}: concepto $\rightarrow$ Prompt $\rightarrow$ herramienta de monitoreo (LangFuse).
    \item \textbf{LightRAG es un esquema de RAG basado en un grafo de conocimiento}, con mayor comprensión semántica y capacidad de actualización incremental que un RAG puramente vectorial.
    \item \textbf{La recuperación de dos niveles} (Low-Level + High-Level) retoma la idea de la recuperación Local y Global de GraphRAG, pero con una implementación más ligera.
    \item \textbf{LangFuse es una herramienta poderosa para entender el mecanismo interno de un Agent}: al hacer trace de todas las llamadas a la API, se puede ver con claridad la entrada y la salida de cada etapa.
    \item \textbf{API principales}: \texttt{ainsert} (indexación incremental) y \texttt{aquery} (recuperación y respuesta), que corresponden a los dos procesos principales, Indexing y Retrieval.
\end{enumerate}

\end{document}
